\section{数值求解方法}

\subsection{守恒量与原始量：变量储存策略}

我的程序同时保存守恒量
$\bm{Q}=[\rho,\rho u,\rho E]^{\mathrm T}$ 和原始量
$\bm{W}=[\rho,u,p]^{\mathrm T}$ 。其中守恒量是时间推进主状态，
原始量由守恒量反解，用于声速、CFL、人工粘性、输出和诊断。

原因讨论如下。

虽然本问题求解的是同一组守恒型方程，但根据推进过程中选取的核心变量不同，
可以采用两种内存储存方案。

\subsubsection{以守恒量 \texorpdfstring{$\bm{Q}$}{Q} 为核心变量}

第一种方案在每个网格点储存
\begin{equation}
\bm{Q}=
\begin{bmatrix}
\rho & \rho u & \rho E
\end{bmatrix}^{\mathrm T},
\end{equation}
并由数值格式直接更新守恒量。MacCormack、有限体积、Roe 和 HLLC
等守恒型数值方法天然以 $\bm{Q}$ 作为推进变量。

该方案的基本流程为
\begin{equation}
\text{内存分配}
\longrightarrow
\bm{Q}
\longrightarrow
\text{时间循环}
\longrightarrow
\text{更新 }\bm{Q}
\longrightarrow
\text{下一时间步}.
\end{equation}

通量、声速、CFL 条件和人工粘性系数通常需要使用原始量，因此每个时间步还需进行
\begin{equation}
\bm{Q}\longrightarrow\bm{W},
\qquad
\bm{W}\longrightarrow\bm{F},
\qquad
\bm{W}\longrightarrow\nu,
\end{equation}
其中 $\nu$ 表示人工粘性系数。控制方程的非线性主要集中在守恒量到原始量的转换、
矢通量构造和人工粘性系数计算中。

\subsubsection{以原始量 \texorpdfstring{$\bm{W}$}{W} 为核心变量}

第二种方案在每个网格点储存
\begin{equation}
\bm{W}=
\begin{bmatrix}
\rho & u & p
\end{bmatrix}^{\mathrm T},
\end{equation}
再由 $\bm{W}$ 组合出守恒量 $\bm{Q}(\bm{W})$、通量
$\bm{F}(\bm{W})$ 和人工粘性系数。其基本流程为
\begin{equation}
\text{内存分配}
\longrightarrow
\bm{W}
\longrightarrow
\bm{Q}(\bm{W}),\ \bm{F}(\bm{W}),\ \nu(\bm{W})
\longrightarrow
\text{更新}
\longrightarrow
\text{下一时间步}.
\end{equation}

Project 0 利用特征关系求解 Riemann 问题的半解析解时，主要采用原始量描述左右状态、
激波、膨胀波和接触间断，因此这种组织方式更符合其理论推导逻辑。

\subsubsection{本程序采用的储存方案}

观察 Euler 方程的守恒形式及 MacCormack 格式可知，时间推进的核心变量天然是
$\bm{Q}$，因此本程序以第一种方案为基础。同时，在内存允许的情况下，为原始量
$\bm{W}$ 也分配长期储存空间，但规定
\begin{equation}
\bm{Q}\text{ 为主状态},
\qquad
\bm{W}\text{ 为由 }\bm{Q}\text{ 同步得到的子状态}.
\end{equation}

$\bm{W}$ 不使用独立的数值格式推进，而是在 $\bm{Q}$ 更新后重新反解。一个完整时间步为
\begin{enumerate}[leftmargin=2em]
    \item 当前主状态 $\bm{Q}^{n}$ 已知；
    \item 由 $\bm{Q}^{n}$ 反解 $\bm{W}^{n}$；
    \item 用 $\bm{W}^{n}$ 计算声速、CFL、人工粘性和通量；
    \item 先对 $\bm{Q}^{n}$ 做人工粘性滤波，得到 $\bar{\bm{Q}}^{n}$；
    \item 执行 MacCormack 预测步，得到 $\widetilde{\bm{Q}}^{n+1}$；
    \item 由 $\widetilde{\bm{Q}}^{n+1}$ 反解预测状态对应的原始量；
    \item 计算预测状态通量 $\overline{\bm{F}}$；
    \item 执行 MacCormack 校正步，得到 $\bm{Q}^{n+1}$；
    \item 由 $\bm{Q}^{n+1}$ 同步得到 $\bm{W}^{n+1}$；
    \item 输出当前状态并进入下一时间步。
\end{enumerate}

这种实现与仅临时计算 $\bm{W}$ 的第一种方案相比，主要区别是
$\bm{W}$ 拥有实际的数组储存空间，因此可直接查询、输出和调试。
代价是需要额外内存，并且必须保证 $\bm{W}$ 始终由最新的 $\bm{Q}$ 同步得到。

\subsection{时间变量储存策略}

\subsubsection{方案 A：全历史存储}

全历史存储为每一个时间层分配空间，保存
\begin{equation}
\bm{Q}^{0},\bm{Q}^{1},\ldots,\bm{Q}^{N_t}.
\end{equation}

该方案可以回看任意时刻、生成动画、进行完整时空分析、检查波传播过程，
并研究误差随时间的演化。但其内存和文件规模随网格数与时间步数的乘积增长，
计算时也会产生较大的缓存压力，因此不适合大型 CFD 计算。

例如，当
\begin{equation}
N_x=100000,\qquad
N_t=100000,\qquad
N_{\mathrm{var}}=3
\end{equation}
且每个 \texttt{double} 占用 $8$ 字节时，仅储存三项守恒量就需要
\begin{equation}
100000\times100000\times3\times8
=240\times10^9\ \text{bytes}
\approx240\ \text{GB}.
\end{equation}

该方案主要适用于小型算例、快速验证、后处理动画、程序调试和保存全部瞬态。

\subsubsection{方案 B：滚动时间层存储}

滚动存储只保存当前计算需要的少数时间层。对于本程序的 MacCormack 格式，
内存中主要保留
\begin{equation}
\bm{Q}^{n},
\qquad
\bar{\bm{Q}}^{n},
\qquad
\widetilde{\bm{Q}}^{n+1},
\qquad
\bm{Q}^{n+1}.
\end{equation}

完成一个时间步后，用 $\bm{Q}^{n+1}$ 覆盖旧的 $\bm{Q}^{n}$，再继续计算。
这种方式内存需求基本不随总时间步数增加，适合实际 CFD 时间推进。

\subsubsection{滚动存储与快照输出}

方案 B 仍可获得方案 A 的中间过程分析能力。本程序在内存中使用滚动存储，
同时每隔 $N$ 个时间步将当前流场写入磁盘快照文件：
\begin{equation}
\text{内存：滚动时间层}
\qquad
\text{磁盘：离散快照}.
\end{equation}

若未来需要研究某个中间状态，可以
\begin{enumerate}[leftmargin=2em]
    \item 找到对应时刻的 snapshot 文件；
    \item 读入其中保存的 $\bm{Q}$ 或 $\bm{W}$；
    \item 将该状态作为新的初始条件；
    \item 从该时间点继续计算。
\end{enumerate}

因此，本程序采用“内存滚动推进、磁盘按需保存快照”的策略，在控制内存消耗的同时，
保留对中间过程进行绘图、分析和重新计算的能力。


\subsection{空间与时间离散}

计算域离散为 $N_x$ 个节点：

\begin{equation}
x_i=x_{\min}+i\Delta x,
\qquad
\Delta x=\frac{x_{\max}-x_{\min}}{N_x-1}.
\end{equation}

当前程序没有采用自适应计算域或自适应网格。主要原因是：本项目需要比较不同
\(\Delta x\)、CFL 数和人工粘性参数下的结果，固定计算域和均匀网格更便于控制变量。

\subsection{CFL 时间步}

\begin{equation}
\Delta t
=
\mathrm{CFL}
\frac{\Delta x}
{\max_i(|u_i|+a_i)},
\qquad
a_i=\sqrt{\frac{\gamma p_i}{\rho_i}}.
\end{equation}

结合 Homework 3 中采用的冯诺依曼稳定性分析，可将离散格式对某一波数扰动的作用
写成放大因子 $G(k)$。在复平面上，稳定性要求任意波数 $k$ 对应的放大因子均满足
\begin{equation}
\left|G(k)\right|\leq 1,
\qquad \forall\,k,
\end{equation}
即所有波数扰动的放大因子轨迹必须位于单位圆内部或单位圆上。CFL 条件正是对
$\Delta t/\Delta x$ 的限制，使离散格式的放大因子不越出复平面上的单位圆，
从而保证任意傅里叶扰动不会随时间步进而无限放大。

\subsection{MacCormack 预测--校正格式}

本文主要采用与课堂公式一致的写法：先进行人工粘性滤波，再使用后向通量差分预测、前向通量差分校正：

\begin{align}
\widetilde{\bm{Q}}_i^{\,n+1}
&=
\bar{\bm{Q}}_i^n
-
\frac{\Delta t}{\Delta x}
\left(
\bm{F}(\bar{\bm{Q}}_i^n)-\bm{F}(\bar{\bm{Q}}_{i-1}^n)
\right),\\
\bm{Q}_i^{n+1}
&=
\frac{1}{2}
\left[
\bar{\bm{Q}}_i^n+\widetilde{\bm{Q}}_i^{\,n+1}
-
\frac{\Delta t}{\Delta x}
\left(
\bm{F}(\widetilde{\bm{Q}}_{i+1}^{\,n+1})
-
\bm{F}(\widetilde{\bm{Q}}_i^{\,n+1})
\right)
\right].
\end{align}

程序也保留了相反方向的写法，即前向预测、后向校正。结果报告中另外选 Sod 算例比较了固定方向和交替方向的影响。

\subsection{人工粘性}

为比较不同物理量对间断位置的识别能力，程序允许在密度 $\rho$、速度 $u$
和压力 $p$ 之间选择人工粘性开关变量。令
\begin{equation}
\phi\in\{\rho,u,p\},
\end{equation}
则统一采用归一化二阶差分构造开关函数：
\begin{equation}
\theta_i(\phi)=
\frac{
\left|\phi_{i+1}-2\phi_i+\phi_{i-1}\right|
}{
\left|\phi_{i+1}-\phi_{i}\right|
+\left|\phi_{i}-\phi_{i-1}\right|
+\varepsilon
},
\end{equation}
其中 $\varepsilon$ 是防止分母为零的小量。人工粘性强度为
\begin{equation}
\nu_i=CFL\cdot(1-CFL)\beta\cdot\theta_i(\phi),
\end{equation}
其中 $\beta$ 为人工粘性经验系数，通常取 \(O(1)\) 量级。CFL 数取
\[
\mathrm{CFL}=\frac{\rho(\mathbf A)\Delta t}{\Delta x},
\qquad
\rho(\mathbf A)=\max_i(|u_i|+c_i).
\]

代码中三种选择共用同一套人工粘性算法，仅切换传入开关函数的
$\phi$ 数组，因此可以在其他参数完全相同的条件下比较
$\rho$、$u$ 和 $p$ 作为开关变量时的数值结果。

\subsection{边界条件}

当前程序使用零梯度外推：

\begin{equation}
\bm{Q}_0=\bm{Q}_1,
\qquad
\bm{Q}_{N_x-1}=\bm{Q}_{N_x-2}.
\end{equation}

这种边界条件适合普通 shock tube 算例。使用时需要保证计算域足够大，
使主要波系在终止时间前不要传播到边界；如果要研究反射边界或周期边界，则需要另写边界处理。
